% Vietnamese translation for OI Public Library
% Source-File: IOI中国国家候选队论文/2003/金恺/金恺.ppt
\documentclass[11pt]{article}
\usepackage{oipl-vn}

\title{Bản trình chiếu: Tối ưu tìm kiếm sâu qua bài chia hình vuông}
\author{Jin Kai}
\date{2003}

\begin{document}
\maketitle

\origtitle{Exploring optimization techniques in depth-first search: from the square dissection problem}
\sourcepath{IOI China candidate team papers / 2003 / Jin Kai / Jin Kai.ppt}

\section{Bài toán chia hình vuông}

Bài toán yêu cầu chia hình vuông lớn gồm \(n\times n\) ô đơn vị thành một số
hình vuông nhỏ hơn, mỗi hình có cạnh nguyên, sao cho số hình vuông nhỏ là ít
nhất. Ràng buộc trong slide là
\[
  1\le n\le 32,
\]
môi trường lập trình là \textit{FreePascal}, bộ nhớ khả dụng
\(64\mathrm{MB}\). Ví dụ \(n=7\) có một phương án tối ưu dùng \(9\) hình
vuông nhỏ.

Nếu \(n\) chẵn, đáp án hiển nhiên không quá \(4\): chia hình vuông thành bốn
hình vuông cạnh \(n/2\). Slide xem đây là trường hợp đơn giản. Khi \(n\) lẻ,
khó nhận ra quy luật toán học trực tiếp, nên trọng tâm chuyển sang tìm kiếm
sâu có cắt tỉa.

\section{Một cận trên bằng quy hoạch động hình chữ nhật}

Trước khi tìm kiếm, slide xét quy hoạch động trên hình chữ nhật. Gọi
\[
  f_{i,j}
\]
là số hình vuông ít nhất để chia hình chữ nhật \(i\times j\). Chuyển trạng
thái tự nhiên là cắt hình chữ nhật bởi một đường thẳng ngang hoặc dọc:
\[
  f_{i,j}=\min\left\{
    \min_k(f_{i,k}+f_{i,j-k}),
    \min_k(f_{k,j}+f_{i-k,j})
  \right\}.
\]

Cách này chỉ cho một cận trên. Với \(n=7\), nó cho kết quả \(10\), không phải
tối ưu, vì phương án tối ưu không nhất thiết bị tách được bởi một đường thẳng
duy nhất thành hai phần độc lập. Tuy vậy giá trị \(f_{n,n}\) thường gần đáp
án tối ưu, nên có thể dùng làm nghiệm ban đầu để giới hạn độ sâu tìm kiếm.

\section{Tìm kiếm ban đầu và cắt tỉa}

Đối tượng tìm kiếm là vị trí và cạnh của từng hình vuông nhỏ. Một thứ tự đơn
giản là đánh số ô \((i,j)\) bằng \((i-1)n+j\), mỗi bước lấy ô chưa phủ có số
nhỏ nhất làm góc trái trên của hình vuông tiếp theo, rồi thử các cạnh khả dĩ.
Cách này đúng nhưng quá chậm.

Các cắt tỉa đầu tiên trong slide gồm:
\begin{itemize}
  \item cắt tỉa đối xứng: quy định cạnh hình vuông ở góc trái trên không lớn
  hơn các hình ở ba góc còn lại, và hình ở góc phải trên không nhỏ hơn hình
  ở góc trái dưới;
  \item cắt tỉa tổng bình phương: nếu diện tích còn lại là \(i\), dùng
  \(Sum_i\) là số ít nhất các số nguyên cần có để tổng bình phương bằng
  \(i\), từ đó lấy làm cận dưới số hình vuông còn phải thêm.
\end{itemize}

Tuy nhiên, với thứ tự quét ô đơn giản, vùng đã phủ thường lồi lõm khó kiểm
soát, làm cận dưới yếu và số nhánh lớn.

\section{Thứ tự tìm kiếm theo đường chia}

Slide đưa ra khái niệm ``hình chữ nhật dạng gần chữ nhật''. Trực giác là biên
giữa phần đã phủ và phần chưa phủ có thể xem như một đường chia đi từ góc
trái dưới đến góc phải trên, chỉ đi sang phải hoặc đi lên. Hình vuông lớn ban
đầu là một trạng thái như vậy; với mỗi trạng thái, có thể chọn một phần lõm
và tách ra một hình vuông mà vẫn giữ dạng này.

Trong tìm kiếm mới, mỗi bước chọn góc trái trên của một phần lõm làm góc
trái trên của hình vuông, và cạnh hình vuông không vượt quá bề rộng phần lõm.
Nhờ đó vùng đã phủ luôn có cấu trúc được kiểm soát. Để tránh đếm lại cùng
một phương án, slide đặt thêm quy tắc thứ tự: hình vuông trước không được nằm
hoàn toàn phía trên hình vuông sau.

Thứ tự mới tạo ra cận dưới mạnh hơn. Nếu trạng thái hiện tại có \(x\) phần
lồi, nó cần ít nhất \(x\) hình vuông nữa, vì ô ở góc phải dưới của mỗi phần
lồi phải thuộc về một hình vuông khác nhau. Thực nghiệm trong slide cho biết
cận này thường chính xác hơn cận \(Sum_i\). Kết hợp thêm đối xứng qua đường
chéo trái trên--phải dưới, thời gian trường hợp chậm \(n=31\) giảm còn
khoảng \(20\) giây trên máy thử nghiệm được ghi trong slide.

\section{Tiền xử lý giá trị trạng thái}

Slide tiếp tục tối ưu bằng ý tưởng biến một số nút trong cây tìm kiếm thành
lá. Nếu đã dùng \(x\) hình vuông và biết vùng chưa phủ cần ít nhất \(y\) hình
vuông nữa, thì mọi lời giải đi qua nút đó có giá trị \(x+y\); không cần mở
rộng nút.

Vấn đề là biểu diễn trạng thái vùng chưa phủ. Một cách trực tiếp dùng bộ
\[
  (a_1,a_2,\ldots,a_n),\qquad a_i\ge a_{i+1},
\]
trong đó \(a_i\) là số ô chưa phủ trên một hàng tương ứng. Cách này tốn
\(n\) byte cho mỗi trạng thái và so sánh hai trạng thái mất \(O(n)\).

Để mã hóa gọn hơn, đặt
\[
  b_i=a_i-a_{i+1}\quad (0\le i\le n),
  \qquad a_0=n,\quad a_{n+1}=0.
\]
Khi đó \(b_i\ge0\) và
\[
  \sum_i b_i=n.
\]
Số bộ \(b\) là số cách đặt \(n\) vật không phân biệt vào \(n+1\) hộp có nhãn,
tức \(\binom{2n}{n}\), với \(n=31\) vẫn không vượt quá \(10^{18}\) theo ghi
chú của slide. Vì \(a\) và \(b\) tương ứng một-một, có thể dùng một số nguyên
kiểu \texttt{comp} để mã hóa trạng thái.

Slide phác thảo bảng
\[
  Tot_{i,j}
\]
là số dãy thỏa \(j\ge a_i\ge a_{i+1}\ge\cdots\ge a_n\ge0\), dùng để chuyển
qua lại giữa bộ \(a\) và số nguyên. Các đoạn mã trong slide chỉ là mẫu chuyển
đổi chỉ số: cộng dồn \(Tot[i,a[i]-1]\) khi mã hóa, và giảm \(a[i]\) theo
\(Tot\) khi giải mã; bản dịch này không chép lại đầy đủ mã khuôn.

\section{BFS trạng thái nhỏ và bảng băm}

Sau khi trạng thái được mã hóa bằng một số nguyên, ta có thể tính trước giá
trị tối ưu của một số trạng thái. Nếu xem mỗi trạng thái là một đỉnh, và thêm
một hình vuông tạo ra một cạnh có độ dài \(1\), thì giá trị \(f_i\) của trạng
thái \(i\) là khoảng cách ngắn nhất từ trạng thái rỗng tới \(i\). Vì mọi cạnh
cùng trọng số, có thể dùng BFS từ trạng thái ban đầu và mở rộng đến khi có
đủ nhiều trạng thái.

Mỗi trạng thái sinh ra được đưa vào bảng băm để tra cứu khi DFS chính đang
chạy. Nếu trạng thái hiện tại có trong bảng, giá trị đã biết cho phép cập
nhật nghiệm tốt nhất và quay lui ngay. Để tránh tra cứu vô ích ở hầu hết các
nút, slide chỉ tính các trạng thái có \(f_i\le c\), với \(c\) là hằng số tự
chọn. Khi đã dùng \(x\) hình vuông và \(x+c+1\) không còn tốt hơn nghiệm
hiện tại, mới cần tra bảng:
\begin{itemize}
  \item nếu không tìm thấy, vùng còn lại cần ít nhất \(c+1\) hình vuông, nên
  nhánh này không thể cải thiện nghiệm;
  \item nếu tìm thấy \(f_i\), cập nhật nghiệm bằng \(x+f_i\) rồi quay lui.
\end{itemize}

Nhờ tiền xử lý này, độ sâu tìm kiếm giảm được khoảng \(c\) tầng. Kết hợp với
các cắt tỉa trước đó, trường hợp \(n=31\) trong bảng thời gian của slide chỉ
còn hơn \(4\) giây.

\section{Tổng kết}

Quá trình giải trong slide gồm bốn bước: dùng DP hình chữ nhật để lấy cận
trên, thêm các cắt tỉa đối xứng và diện tích, đổi thứ tự tìm kiếm theo cấu
trúc đường chia, rồi mã hóa trạng thái để tiền xử lý một phần cây tìm kiếm.

Thông điệp chung là thứ tự và đối tượng tìm kiếm là nền tảng của hiệu quả,
nhưng cắt tỉa đơn thuần không làm thay đổi bản chất độ phức tạp. Khi bộ nhớ
cho phép, có thể hy sinh không gian để tính trước giá trị của các trạng thái
nhỏ, từ đó giảm độ sâu DFS. Bước khó nhất của tối ưu này là biểu diễn trạng
thái đủ gọn để tra cứu nhanh.

\end{document}
